% !TEX root = ../main.tex
\begin{ouverture}[photons]{25}{Photons et effet photoélectrique}{Partie 4 · Ondes et signaux}
{Photon\enspace·\enspace Effet photoélectrique\enspace·\enspace Travail d'extraction\enspace·\enspace Cellule photovoltaïque}
\objectif{Décrire l'effet photoélectrique, ses caractéristiques et son importance historique.}
\objectif{Interpréter qualitativement l'effet photoélectrique à l'aide du modèle particulaire de la lumière.}
\objectif{Établir, par un bilan d'énergie, la relation entre l'énergie cinétique des électrons et la fréquence.}
\objectif{Expliquer le fonctionnement d'une cellule photoélectrique et citer des applications de l'interaction
photon-matière.}
\objectif{Déterminer le rendement d'une cellule photovoltaïque, à l'aide d'un programme Python.}
\prerequis{Photon, énergie $E=h\nu$, niveaux d'énergie et spectres (classe de première).}
\prerequis{Électronvolt, énergie cinétique (classe de première, chapitre~15).}
\prerequis{Puissance électrique $P=UI$, rendement d'un convertisseur (classe de première).}
\end{ouverture}
\essentiel
\essSection{1}{Le photon}
\begin{definition}[Le modèle particulaire de la lumière]
La lumière peut être décrite comme un flux de particules, les \textbf{photons}. Un photon a une masse nulle, se déplace à la
vitesse de la lumière $c$ dans le vide et transporte l'énergie
\[E=h\nu=\frac{h\,c}{\lambda},\]
où $\nu$ est la fréquence de l'onde associée, $\lambda$ sa longueur d'onde et $h=\qty{6,63e-34}{J.s}$ la constante de Planck.
On exprime souvent ces énergies en électronvolts : $\qty{1}{eV}=\qty{1,60e-19}{J}$.
\end{definition}
\begin{exemple}
Un photon de lumière verte ($\lambda=\qty{500}{nm}$) transporte $E=\dfrac{\num{6.63e-34}\times\num{3.00e8}}{\num{500e-9}}=\qty{3,98e-19}{J}$, soit
\qty{2,49}{eV}. Un pointeur laser de \qty{1}{mW} émettant à \qty{633}{nm} envoie donc $\num{3.2e15}$ photons par seconde.
\end{exemple}
\begin{propriete}[Absorption et émission de photons]
Les énergies d'un atome, d'une molécule ou d'un solide sont quantifiées. Un photon est absorbé si son énergie permet une
transition vers un niveau d'énergie supérieur ; inversement, la désexcitation s'accompagne de l'émission d'un photon
d'énergie égale à l'écart entre les niveaux. L'interaction photon-matière est à la base des spectroscopies UV-visible et
infrarouge, des diodes électroluminescentes, des lasers, des capteurs d'images et des cellules photovoltaïques.
\end{propriete}
\essSection{2}{L'effet photoélectrique}
\begin{definition}[Effet photoélectrique]
L'\textbf{effet photoélectrique} est l'émission d'électrons par un métal éclairé par une lumière de fréquence
suffisamment élevée. Observé par Heinrich Hertz en 1887, il a été expliqué en 1905 par Albert Einstein grâce au modèle
particulaire de la lumière, ce qui lui valut le prix Nobel de physique en 1921 ; Robert Millikan en a vérifié les
prévisions en 1916.
\end{definition}
\begin{propriete}[Caractéristiques]
\begin{itemize}
\item Il existe une \textbf{fréquence seuil} $\nu_0$, propre à chaque métal : en dessous, aucun électron n'est émis, quelle
que soit l'intensité de la lumière.
\item Au-dessus du seuil, l'émission est immédiate, même pour une lumière très faible.
\item L'énergie cinétique maximale des électrons dépend de la fréquence de la lumière, pas de son intensité ; le
nombre d'électrons émis par seconde est proportionnel à l'intensité.
\end{itemize}
Le modèle ondulatoire, où l'énergie de la lumière s'accumule progressivement, ne peut expliquer ni le seuil ni
l'émission immédiate.
\end{propriete}
\schema{25-effet}
\begin{savoirfaire}[Établir la relation d'Einstein par un bilan d'énergie]
Un photon absorbé cède toute son énergie à un seul électron. Il faut fournir au moins le \textbf{travail d'extraction}
$W_\mathrm{ext}$ pour arracher un électron au métal ; le reste de l'énergie devient énergie cinétique. Pour les électrons
les moins liés :
\[h\nu=W_\mathrm{ext}+E_\mathrm{c,max}\qquad\text{soit}\qquad E_\mathrm{c,max}=h\nu-W_\mathrm{ext}.\]
L'émission n'a lieu que si $h\nu\geqslant W_\mathrm{ext}$, d'où la fréquence seuil $\nu_0=\dfrac{W_\mathrm{ext}}{h}$ et la longueur d'onde seuil
$\lambda_0=\dfrac{h\,c}{W_\mathrm{ext}}$.
\end{savoirfaire}
\schema{25-ec}
\begin{exemple}
Pour le zinc, $W_\mathrm{ext}=\qty{4,3}{eV}$ : $\lambda_0=\dfrac{h\,c}{W_\mathrm{ext}}=\qty{289}{nm}$. Seul le rayonnement ultraviolet arrache des
électrons au zinc ; une lumière visible, même intense, n'a aucun effet. Pour le sodium ($W_\mathrm{ext}=\qty{2,28}{eV}$),
$\nu_0=\qty{5,5e14}{Hz}$, soit $\lambda_0=\qty{545}{nm}$ : la lumière bleue ou verte suffit.
\end{exemple}
\essSection{3}{Cellules photoélectriques et photovoltaïques}
\begin{propriete}[Convertir la lumière en courant]
Dans une \textbf{cellule photoélectrique}, une cathode métallique éclairée émet des électrons recueillis par une anode :
l'intensité du courant, proportionnelle au nombre de photons reçus, mesure l'éclairement. Dans une \textbf{cellule
photovoltaïque}, un photon d'énergie supérieure à un seuil, le \textbf{gap} du semi-conducteur (\qty{1,12}{eV} pour le
silicium, soit $\lambda<\qty{1,11}{\micro m}$), libère un électron que la structure de la cellule met en mouvement : la cellule
fournit une puissance électrique. Son \textbf{rendement} est le rapport
\[\eta=\frac{P_\text{électrique}}{P_\text{lumineuse reçue}},\qquad P_\text{lumineuse reçue}=E\times S,\]
où $E$ est l'éclairement énergétique (en \unit{W.m^{-2}}) et $S$ la surface de la cellule. Il atteint 15 à 23~\% pour les
panneaux du commerce.
\end{propriete}
\schema{25-cellule}
\begin{python}
U = [0, 1.0, 2.0, 3.0, 4.0, 4.5, 5.0, 5.5, 6.0]           # V
I = [0.85, 0.85, 0.84, 0.83, 0.80, 0.76, 0.66, 0.42, 0]   # A
P = [u * i for u, i in zip(U, I)]                          # W
k = P.index(max(P))                                        # puissance maximale
print(U[k], round(P[k], 2), round(100 * P[k] / (1000 * 0.025), 1))
\end{python}
\begin{console}
4.5 3.42 13.7
\end{console}
\noindent\capacite{Capacité numérique}Les listes contiennent les mesures de la tension et de l'intensité délivrées par un petit
panneau de \qty{0,025}{m^2} sous un éclairement de \qty{1000}{W.m^{-2}}. Le programme calcule la puissance $P=UI$ pour chaque
mesure, repère son maximum, \qty{3,42}{W} sous \qty{4,5}{V}, et en déduit le rendement : 13,7~\%.
\begin{attention}
L'intensité d'une lumière fixe le nombre de photons par seconde, sa fréquence l'énergie de chaque photon. Augmenter
l'intensité d'une lumière de fréquence inférieure au seuil n'arrache aucun électron. Penser à convertir les énergies :
$\qty{1}{eV}=\qty{1,60e-19}{J}$.
\end{attention}
\begin{experience}[Déterminer le rendement d'une cellule photovoltaïque]
On éclaire une cellule photovoltaïque avec une lampe et l'on mesure l'éclairement reçu avec un capteur. La cellule
alimente une résistance variable : pour chaque réglage, on relève la tension à ses bornes et l'intensité du courant. On
trace $I$ et $P=UI$ en fonction de $U$, on repère la puissance maximale et l'on calcule le rendement.
\end{experience}
\begin{rediger}
Écrire le bilan d'énergie en mots avant de l'écrire en équation : « énergie du photon = travail d'extraction + énergie
cinétique de l'électron ». Préciser l'unité de chaque énergie et ne mélanger joules et électronvolts qu'après conversion.
\end{rediger}
\exercices

\rubrique{Découverte}
\begin{blocLivre}{ex-25-1}
\begin{exercice}[1]{L'énergie d'un photon}
\begin{donnees}
$h=\qty{6,63e-34}{J.s}$ ; $c=\qty{3,00e8}{m.s^{-1}}$ ; $\qty{1}{eV}=\qty{1,60e-19}{J}$.
\end{donnees}
\begin{questions}
\item Calculer, en électronvolts, l'énergie d'un photon de lumière violette (\qty{400}{nm}), puis rouge (\qty{700}{nm}).
\item Même question pour un photon ultraviolet de \qty{254}{nm} (lampe germicide) et un photon infrarouge de \qty{1000}{nm}.
\item Rompre une liaison carbone-carbone demande \qty{3,6}{eV}. Quels photons peuvent la rompre ? Calculer la longueur
d'onde maximale correspondante.
\item Pourquoi les ultraviolets sont-ils dangereux pour la peau, alors que l'infrarouge ne fait que la chauffer ?
\end{questions}
\end{exercice}
\end{blocLivre}
\begin{blocLivre}{ex-25-2}
\begin{exercice}[1]{Les seuils photoélectriques}
\begin{center}
\begin{tabular}{lcccc}
\toprule
Métal & césium & sodium & zinc & cuivre \\
$W_\mathrm{ext}$ (eV) & 2,1 & 2,28 & 4,3 & 4,7 \\
\bottomrule
\end{tabular}
\end{center}
\begin{donnees}
$h=\qty{6,63e-34}{J.s}$ ; $c=\qty{3,00e8}{m.s^{-1}}$ ; $\qty{1}{eV}=\qty{1,60e-19}{J}$.
\end{donnees}
\begin{questions}
\item Définir le travail d'extraction d'un métal.
\item Calculer la longueur d'onde seuil de chaque métal.
\item Quels métaux émettent des électrons lorsqu'ils sont éclairés en lumière visible (400 à \qty{800}{nm}) ?
\item Pourquoi utilise-t-on du césium dans les cellules photoélectriques sensibles à la lumière visible ?
\end{questions}
\end{exercice}
\end{blocLivre}
\begin{blocLivre}{ex-25-3}
\begin{exercice}[1]{Des électrons arrachés au sodium}
Une plaque de sodium ($W_\mathrm{ext}=\qty{2,28}{eV}$) est éclairée par une lumière de longueur d'onde \qty{400}{nm}.
\begin{donnees}
$h=\qty{6,63e-34}{J.s}$ ; $c=\qty{3,00e8}{m.s^{-1}}$ ; $\qty{1}{eV}=\qty{1,60e-19}{J}$ ; masse de l'électron : \qty{9,11e-31}{kg}.
\end{donnees}
\begin{questions}
\item Calculer l'énergie d'un photon incident, en électronvolts.
\item Écrire le bilan d'énergie et calculer l'énergie cinétique maximale des électrons émis, en eV puis en joules.
\item En déduire leur vitesse maximale.
\item Que se passe-t-il avec une lumière de \qty{600}{nm} ? Et si l'on double l'intensité de cette lumière ?
\end{questions}
\end{exercice}
\end{blocLivre}
\begin{blocLivre}{ex-25-4}
\begin{exercice}[1]{Onde ou particule ?}
\begin{questions}
\item Dans le modèle ondulatoire, l'énergie apportée par la lumière est proportionnelle à son intensité et s'accumule au
cours du temps. Que prévoit ce modèle si l'on éclaire longtemps un métal avec une lumière faible de basse fréquence ?
\item Qu'observe-t-on en réalité ?
\item Comment le modèle particulaire explique-t-il l'existence d'une fréquence seuil ?
\item La lumière est-elle une onde ou un flux de particules ? Citer un phénomène qui illustre chaque aspect.
\end{questions}
\end{exercice}
\end{blocLivre}
\begin{blocLivre}{ex-25-5}
\begin{exercice}[1]{Combien de photons ?}
\begin{donnees}
$h=\qty{6,63e-34}{J.s}$ ; $c=\qty{3,00e8}{m.s^{-1}}$.
\end{donnees}
\begin{questions}
\item Calculer l'énergie d'un photon émis par un pointeur laser vert de longueur d'onde \qty{532}{nm}.
\item Le pointeur a une puissance de \qty{1,0}{mW}. Combien de photons émet-il par seconde ?
\item En plein soleil, un mètre carré reçoit environ \qty{1000}{W}. En prenant une longueur d'onde moyenne de
\qty{550}{nm}, estimer le nombre de photons reçus par seconde.
\item Pourquoi ne perçoit-on pas le caractère « granulaire » de la lumière dans la vie courante ?
\end{questions}
\end{exercice}
\end{blocLivre}
\rubrique{Entraînement}
\begin{blocLivre}{ex-25-6}
\begin{exercice}[2]{À la manière de Millikan}
On éclaire une cathode de sodium avec des lumières de fréquences différentes. Pour chaque fréquence, on mesure la tension
d'arrêt $U_\mathrm{a}$, la plus petite tension qui empêche les électrons les plus rapides d'atteindre l'anode : $E_\mathrm{c,max}=e\,U_\mathrm{a}$.
\begin{center}
\begin{tabular}{lcccc}
\toprule
$\nu$ ($10^{14}$ Hz) & 6,9 & 7,4 & 8,2 & 9,6 \\
$U_\mathrm{a}$ (V) & 0,57 & 0,78 & 1,12 & 1,69 \\
\bottomrule
\end{tabular}
\end{center}
\begin{donnees}
$e=\qty{1,60e-19}{C}$.
\end{donnees}
\begin{questions}
\item Montrer que la relation d'Einstein s'écrit $U_\mathrm{a}=\dfrac{h}{e}\,\nu-\dfrac{W_\mathrm{ext}}{e}$.
\item À l'aide des mesures extrêmes, déterminer le coefficient directeur et l'ordonnée à l'origine de la droite.
\item En déduire la constante de Planck et le travail d'extraction du sodium, en eV.
\item Calculer la fréquence seuil. Pourquoi l'expérience de Millikan a-t-elle convaincu les physiciens ?
\end{questions}
\end{exercice}
\end{blocLivre}
\begin{blocLivre}{ex-25-7}
\begin{exercice}[2]{L'électroscope de Hallwachs}
Une plaque de zinc, posée sur un électroscope, est chargée négativement : les feuilles de l'électroscope s'écartent.
\begin{questions}
\item On éclaire la plaque avec une lampe à ultraviolets : les feuilles retombent. Interpréter.
\item Avec une lampe à lumière visible très puissante, rien ne se passe. Expliquer, sachant que $W_\mathrm{ext}=\qty{4,3}{eV}$
pour le zinc.
\item La plaque est maintenant chargée positivement. Les ultraviolets la déchargent-ils ? Justifier.
\item On interpose une plaque de verre, qui absorbe les ultraviolets, entre la lampe et le zinc. Que se passe-t-il ?
\end{questions}
\end{exercice}
\end{blocLivre}
\begin{blocLivre}{ex-25-8}
\begin{exercice}[2]{Le rendement d'un panneau en Python}
On reprend le programme de l'essentiel.
\begin{questions}
\item \capacite{Capacité numérique}Que contient la liste \code{P} ? Que représente l'indice \code{k} ?
\item Justifier l'expression du rendement calculée dans la dernière ligne.
\item Le panneau, exposé sous un éclairement de \qty{800}{W.m^{-2}}, délivre les mêmes tensions et des intensités réduites
de 20~\%. Quelles lignes faut-il modifier ? Le rendement change-t-il ?
\item Pourquoi un panneau relié directement à un appareil quelconque ne fonctionne-t-il pas forcément à sa puissance
maximale ?
\end{questions}
\end{exercice}
\end{blocLivre}
\begin{blocLivre}{ex-25-9}
\begin{exercice}[2]{La diode électroluminescente}
Une LED rouge émet une lumière de longueur d'onde \qty{630}{nm}. Elle consomme une puissance électrique de \qty{1,0}{W} et émet
une puissance lumineuse de \qty{0,35}{W}.
\begin{donnees}
$h=\qty{6,63e-34}{J.s}$ ; $c=\qty{3,00e8}{m.s^{-1}}$ ; $\qty{1}{eV}=\qty{1,60e-19}{J}$.
\end{donnees}
\begin{questions}
\item La LED fonctionne-t-elle par absorption ou par émission de photons ? Calculer l'énergie des photons émis, en eV.
\item Une LED ne s'allume qu'à partir d'une tension de seuil voisine de l'énergie de ses photons divisée par $e$. Estimer
cette tension pour la LED rouge.
\item Calculer le nombre de photons émis par seconde.
\item Calculer le rendement de la LED. Que devient le reste de l'énergie ?
\end{questions}
\end{exercice}
\end{blocLivre}
\begin{blocLivre}{ex-25-10}
\begin{exercice}[2]{Les limites d'une cellule au silicium}
Dans une cellule au silicium, un photon n'est absorbé que si son énergie dépasse le gap, \qty{1,12}{eV} ; l'énergie en
excès est perdue sous forme d'énergie interne.
\begin{donnees}
$h=\qty{6,63e-34}{J.s}$ ; $c=\qty{3,00e8}{m.s^{-1}}$ ; $\qty{1}{eV}=\qty{1,60e-19}{J}$.
\end{donnees}
\begin{questions}
\item Calculer la longueur d'onde maximale des photons absorbés.
\item Un photon de \qty{1500}{nm} est-il absorbé ? Justifier.
\item Pour un photon de \qty{400}{nm}, calculer l'énergie en excès et la fraction de son énergie perdue.
\item Expliquer pourquoi le rendement d'une cellule au silicium ne peut pas approcher 100~\%.
\end{questions}
\end{exercice}
\end{blocLivre}
\begin{blocLivre}{ex-25-11}
\begin{exercice}[2]{Le capteur d'un appareil photo}
Chaque pixel du capteur d'un appareil photo est un carré de \qty{4}{\micro m} de côté. Un photon reçu libère un électron avec
une probabilité de 60~\% ; le pixel sature au-delà de \num{30000} électrons. Pendant une pose de \qty{0,010}{s}, le pixel reçoit
une lumière verte de \qty{550}{nm} sous un éclairement de \qty{0,10}{W.m^{-2}}.
\begin{donnees}
$h=\qty{6,63e-34}{J.s}$ ; $c=\qty{3,00e8}{m.s^{-1}}$.
\end{donnees}
\begin{questions}
\item Calculer l'énergie lumineuse reçue par le pixel pendant la pose.
\item En déduire le nombre de photons reçus, puis le nombre d'électrons libérés.
\item Le pixel est-il saturé ? Au bout de quelle durée de pose le serait-il ?
\item Quel phénomène, étudié dans ce chapitre, est à la base du fonctionnement de ce capteur ?
\end{questions}
\end{exercice}
\end{blocLivre}
\rubrique{Objectif bac}
\begin{blocLivre}{ex-25-12}
\begin{exercice}[3]{Des panneaux solaires sur le toit}
\begin{docu}{1}{Le panneau}
Un panneau photovoltaïque de \qty{1,7}{m^2} a une puissance « crête » de \qty{400}{W} : c'est la puissance qu'il fournit sous un
éclairement de \qty{1000}{W.m^{-2}}, à \qty{25}{\degreeCelsius}.
\end{docu}
\begin{docu}{2}{L'installation}
Une maison est équipée de 10 panneaux orientés au sud. Dans sa région, un mètre carré reçoit chaque année
\qty{1300}{kWh} d'énergie solaire. Les pertes (température, câbles, onduleur) réduisent la production de 20~\%. Un foyer
consomme en moyenne \qty{4500}{kWh} d'électricité par an.
\end{docu}
\begin{questions}
\item Calculer le rendement d'un panneau.
\item Calculer la puissance crête de l'installation.
\item \initiative Estimer l'énergie électrique produite en un an. Couvre-t-elle la consommation du foyer ?
\item Pourquoi la production réelle varie-t-elle beaucoup au cours de la journée et de l'année ?
\end{questions}
\end{exercice}
\end{blocLivre}
\begin{blocLivre}{ex-25-13}
\begin{exercice}[3]{Détecter un photon unique}
\begin{docu}{1}{Le photomultiplicateur}
Un photomultiplicateur est un tube à vide dont la photocathode, de travail d'extraction \qty{2,0}{eV}, émet un électron
lorsqu'elle absorbe un photon. Cet électron est accéléré vers une série d'électrodes qui multiplient le nombre
d'électrons : au total, chaque photoélectron produit une impulsion de $10^6$ électrons.
\end{docu}
\begin{docu}{2}{Ses performances}
Un photon de longueur d'onde utile produit un photoélectron avec une probabilité de 25~\%.
\end{docu}
\begin{donnees}
$h=\qty{6,63e-34}{J.s}$ ; $c=\qty{3,00e8}{m.s^{-1}}$ ; $e=\qty{1,60e-19}{C}$.
\end{donnees}
\begin{questions}
\item Calculer la longueur d'onde seuil de la photocathode.
\item Un photon de \qty{420}{nm} atteint la photocathode. Calculer l'énergie cinétique maximale de l'électron émis.
\item \initiative Calculer la charge électrique d'une impulsion. Sur \num{1000} photons reçus, combien d'impulsions compte-t-on en
moyenne ?
\item Pourquoi cet instrument a-t-il besoin de l'effet photoélectrique, et pas seulement de l'effet thermique de la
lumière ?
\end{questions}
\end{exercice}
\end{blocLivre}
\begin{blocLivre}{ex-25-14}
\begin{exercice}[3]{La cellule photoélectrique au césium}
\begin{docu}{1}{La cellule}
Une cellule photoélectrique comporte une cathode de césium ($W_\mathrm{ext}=\qty{2,1}{eV}$) et une anode qui recueille les
électrons émis. Elle reçoit une lumière bleue de \qty{450}{nm}, de puissance \qty{1,0}{\micro W}.
\end{docu}
\begin{docu}{2}{Le rendement quantique}
Seuls 20~\% des photons reçus libèrent un électron qui atteint l'anode.
\end{docu}
\begin{donnees}
$h=\qty{6,63e-34}{J.s}$ ; $c=\qty{3,00e8}{m.s^{-1}}$ ; $e=\qty{1,60e-19}{C}$.
\end{donnees}
\begin{questions}
\item Calculer le nombre de photons reçus par seconde.
\item \initiative Calculer l'intensité du courant produit par la cellule.
\item Calculer l'énergie cinétique maximale des électrons émis, en eV.
\item Comment évoluent l'intensité du courant et l'énergie cinétique maximale si l'on double la puissance lumineuse ? si
l'on remplace la lumière bleue par une lumière rouge de \qty{700}{nm} ?
\end{questions}
\end{exercice}
\end{blocLivre}
\begin{blocLivre}{ex-25-15}
\begin{exercice}[3]{La couleur du laser hélium-néon}
\begin{docu}{1}{Un laser historique}
Dans un laser hélium-néon, des décharges électriques portent les atomes de néon dans un niveau excité de \qty{20,66}{eV}.
L'émission laser correspond à leur désexcitation vers le niveau de \qty{18,70}{eV}.
\end{docu}
\schema{25-ex15}
\begin{donnees}
$h=\qty{6,63e-34}{J.s}$ ; $c=\qty{3,00e8}{m.s^{-1}}$ ; $\qty{1}{eV}=\qty{1,60e-19}{J}$.
\end{donnees}
\begin{questions}
\item S'agit-il d'une absorption ou d'une émission de photon ? Calculer l'énergie du photon émis.
\item \initiative Calculer la longueur d'onde de la lumière laser. Quelle est sa couleur ?
\item La valeur tabulée est \qty{632,8}{nm}. Avec $hc=\qty{1239,84}{eV.nm}$, recalculer la longueur d'onde et commenter.
\item Pourquoi la lumière d'un laser est-elle presque monochromatique, contrairement à celle d'une lampe à
incandescence ?
\end{questions}
\end{exercice}
\end{blocLivre}
